\usepackage{amsmath}
\usepackage{mathtools}
\usepackage{tikz}
%\usepackage{movie15}
\usepackage{algorithm} % must read after hyperref
\usepackage{algorithmic}
\usepackage{subfigure,MnSymbol}
\usepackage{graphicx,color,colortbl}
%\usepackage[footnotesize,bf]{caption}
\usepackage{bm,amsfonts,graphics}
\usepackage{tikz,tkz-graph,tkz-arith,tkz-berge}
\usetikzlibrary{arrows,shapes,positioning}

\graphicspath{{./}{figures/}}
\usepackage{multicol}

\newtheorem{thm}{Theorem}
\newtheorem{lem}{Lemma}
\newtheorem{prop}{Proposition}
\newtheorem{rem}{Remark}

\newcommand{\jcols}[4]{
\begin{columns}[T]\begin{column}[T]{#3\textwidth} #1 \end{column}
\begin{column}[T]{#4\textwidth} #2 \end{column}
\end{columns}}

\newcommand{\jcolsb}[4]
{\begin{columns}[b]
\begin{column}{#3\textwidth} #1 \vspace{0pt}
\end{column}
\begin{column}{#4\textwidth} #2 \vspace{0pt}
\end{column}
\end{columns}}

\newcommand{\jcolsc}[4]
{\begin{columns}[c]
\begin{column}{#3\textwidth} #1 \vspace{0pt}
\end{column}
\begin{column}{#4\textwidth} #2 \vspace{0pt}
\end{column}
\end{columns}}
\mode<presentation>
{
	\usetheme{ACL}
        \setbeamercovered{transparent=0}
}

\makeatletter
\newdimen\labelwidthi
\newdimen\labelwidthii
\newdimen\labelwidthiii
\newdimen\labelwidthiv
\def\normal@labelsep{0.2em}
\labelsep\normal@labelsep
\settowidth{\labelwidthi}{iii}
\settowidth{\labelwidthii}{ii}
\settowidth{\labelwidthiii}{iii}
%\settowidth{\labelwidthiv}{i}
\setlength\leftmargini{0pt}
\setlength\leftmarginii{1.5ex}
\setlength\leftmarginiii{1.5ex}

\leftmargini\labelwidthi    \advance\leftmargini\labelsep
\leftmarginii\labelwidthii  \advance\leftmarginii\labelsep
\leftmarginii\labelwidthiii \advance\leftmarginiii\labelsep

%\leftmarginii\labelwidthiv  \advance\leftmarginiv\labelsep
%\def\setleftmargin#1#2{\settowidth{\@tempdima}{#2}\labelsep\normal@labelsep
%  \csname labelwidth#1\endcsname\@tempdima
%  \@tempdimb\@tempdima \advance\@tempdimb\labelsep
%  \csname leftmargin#1\endcsname\@tempdimb}
%\def\@listI{\leftmargin\leftmargini
%  \labelwidth\labelwidthi \labelsep\normal@labelsep
%  \topsep \z@ \partopsep\z@ \parsep\z@ \itemsep\z@
%  \listparindent 1em}
%\def\@listii{\leftmargin\leftmarginii
%  \labelwidth\labelwidthii \labelsep\normal@labelsep
%  \topsep\z@ \partopsep\z@ \parsep\z@ \itemsep\z@
%  \listparindent 1em}
%\def\@listiii{\leftmargin\leftmarginiii
%  \labelwidth\labelwidthiii \labelsep\normal@labelsep
%  \topsep\z@ \partopsep\z@ \parsep\z@ \itemsep\z@
%  \listparindent 1em}
%\def\@listiv{\leftmargin\leftmarginiv
%  \labelwidth\labelwidthiv \labelsep\normal@labelsep
%  \topsep\z@ \partopsep\z@ \parsep\z@ \itemsep\z@
%  \listparindent 1em}
%\let\@listi\@listI
%\@listi
\makeatother

\logo{
\includegraphics[height = 0.6cm,trim=0 50 0 0,clip]{AO-logo-high_color-top-MIT}
\includegraphics[height = 0.5cm,trim=0 0 0 0,clip]{ACL_logo}
}

\setbeamerfont{frametitle}{size=\large,series=\bfseries}
%\setbeamertemplate{frametitle}[default][center]
\setbeamerfont{title}{size=\large,series=\bfseries}
\setbeamerfont{framesubtitle}{series=\mdseries,size=\large}
\graphicspath{{./}{figures/}}
\usepackage[square, numbers, comma]{natbib}	% use Natbib for author-name
									% citations. In a presentation,
									% this is more legible than pure
									% numbers, like [21]. Google
									% "natbib" for more information.
									% Instead of "\cite{abc}", use
									% "\citep{abc}" or "\citet{abc}"
\bibliographystyle{plainnat}		% Natbib specific
%\bibpunct{(}{)}{,}{a}{}{,}			% Natbib specific

\definecolor{Maroon}{RGB}{112,0,0}
\definecolor{Crimson}{RGB}{200,0,0}
\definecolor{purple}{RGB}{112,0,112}

\newcommand{\bee}{\begin{enumerate}}
\newcommand{\eee}{\end{enumerate}}
\newcommand{\bi}{\begin{itemize}}
\newcommand{\bii}{\begin{itemize}\small}
\newcommand{\beee}{\begin{enumerate}\scriptsize}
\newcommand{\biii}{\begin{itemize}\footnotesize}
\newcommand{\ei}{\end{itemize}}
\newcommand{\jtem}{\vfill\item}
\newcommand{\bea}{\begin{eqnarray}}
\newcommand{\eea}{\end{eqnarray}}
\newcommand{\argmax}{\operatornamewithlimits{argmax}}
\newcommand{\njra}{\color{red} $\Rightarrow$~\color{black}}
\newcommand{\nred}[1]{{\tiny \color{red} XX #1 XX\color{black}}}
%\newcommand{\jframe}[1]{\begin{frame}\nred{#1}\end{frame}}
\newcommand{\jframe}[2]{\begin{frame}\frametitle{#1} #2 \end{frame}}
\newcommand{\jframeb}[2]{\begin{frame}[t]{#1} \begin{itemize} #2 \end{itemize} \end{frame}}
\newcommand{\jframet}[2]{\begin{frame}[t]\frametitle{#1} #2 \end{frame}}
\newcommand{\ji}[2]{\begin{itemize}\setlength{\itemsep}{#1mm}{#2}\end{itemize}}
\newcommand{\jen}[2]{\begin{enumerate}\setlength{\itemsep}{#1mm}{#2}\end{enumerate}}
\newcommand{\bc}{\begin{center}}
\newcommand{\ec}{\end{center}}
\newcommand{\todo}[1]{{\color{magenta} XX- #1 -XX}}
\newcommand{\XX}[1]{{\color{red} XX- #1 -XX}}
\newcommand{\alertr}[1]{\textbf{\color{red} #1}}

\setbeamercovered{transparent=10}
\setbeamercolor{title}{bg = black!10}
\setbeamercolor{frametitle}{bg = black!0}
\setbeamercolor{thm}{bg = black!0}
%\setbeamercolor{alerted text}{fg=blue}
%\setbeamerfont{alerted text}{series = \bfseries}

%{use = palette secondary, bg = palette secondary.bg}

\setbeamercolor{normal text}{fg=black}  
\setbeamertemplate{enumerate items}[ball]
\setbeamertemplate{enumerate subitem}{(\alph{enumii})}
\setbeamertemplate{enumerate subsubitems}[triangle]
\setbeamertemplate{itemize item}{$\filledmedtriangleright$}
\setbeamertemplate{itemize subitem}[circle]
\setbeamertemplate{itemize subitem}{\tiny $\stackrel{\color{red}\filledmedsquare\color{black}}{\hbox{\vrule width 0ex height .25ex}}$}

\setbeamertemplate{navigation symbols}{}
\setbeamersize{text margin left=5mm, text margin right=8mm}

\tikzstyle{input} = [coordinate] %
\tikzstyle{block} = [draw, fill=red!20, rectangle, minimum height=3em, minimum width=6em] %
\tikzstyle{block2} = [draw, fill=blue!20, diamond, minimum height=3em, minimum width=6em] %
